\documentclass[a4paper]{article}     
\usepackage{amsfonts}
\usepackage{amsmath}
\usepackage{amssymb}
\usepackage{graphicx}

\begin{document}

\noindent \large Auteur: Xander De Blauwer \\
\noindent \large Studierichting: Informatica\\
\noindent \large Oefeningenreeks 1C nr. 12.4\\

\medskip

\normalsize





$z_1 = \sqrt{2} + 0i$ \\

$z_2 = -1 + i$ \\

Modulus: \\

$|z_1| = \sqrt{2}$ \\

$|z_2| = \sqrt{2}$ \\


Argumenten: \\

$\tan\theta_1 = 0$ \\

$\theta_1 = \operatorname{Bgtan}(0) = 0$ \\ 

$\tan\theta_2 = -1$ \\

$\theta_2 = \operatorname{Bgtan}(-1) = -\dfrac{\pi}{4}$ \\

$\theta_2 = -\dfrac{\pi}{4} + \pi = \dfrac{3\pi}{4}$ \\



$\arg\left(\dfrac{z_1}{z_2}\right) = \theta_1 - \theta_2$ \\

$\arg\left(\dfrac{z_1}{z_2}\right) = 0 - \dfrac{3\pi}{4} = -\dfrac{3\pi}{4}$ \\ 


Goniometrische vorm: \\

$\dfrac{z_1}{z_2} = \cos\left(-\dfrac{3\pi}{4}\right) + i \sin\left(-\dfrac{3\pi}{4}\right)$ \\

$\dfrac{z_1}{z_2} = \operatorname{cis}\left(-\dfrac{3\pi}{4}\right)$ \\


\begin{thebibliography}{999}
\end{thebibliography}
\end{document} 
